% ============================================================================
% 01_capitulo1_complementos.tex
% Complementos e substituicoes para o Capitulo 1 -- INTRODUCAO.
%
% Como usar:
%   (A) Inserir o paragrafo do bloco 1.A no final da Justificativa (Cap. 1.1),
%       como ancora cientifica recente (Aragao 2018 + Silva-Junior 2025).
%   (B) Substituir os 6 objetivos especificos da Secao 1.2.2 pelo bloco 1.B.
%   (C) Substituir a Estrutura da Monografia (Secao 1.3) pelo bloco 1.C, que
%       descreve os novos Capitulos 5 e 6.
% ============================================================================

% ============================================================================
% (A) PARAGRAFO DE ANCORA CIENTIFICA RECENTE -- inserir ao FIM da Secao 1.1
%     ("Justificativa"), logo antes do paragrafo final que comeca em
%     "Dada a magnitude e complexidade das queimadas...".
% ============================================================================

A urgência do tema é corroborada por publicações recentes na literatura internacional. Aragão \emph{et al.}~\cite{aragao2018amazon} documentaram, no periódico \emph{Nature Communications}, que desde o início do século~XXI o fogo associado à seca passou a \emph{neutralizar} os ganhos obtidos com a redução do desmatamento na Amazônia. Mais recentemente, Silva-Junior \emph{et al.}~\cite{silvajunior2025amazon}, em \emph{Biogeosciences}, mostraram que 2024 marcou a pior degradação por fogo em mais de duas décadas: \SI{3,3}{\mega\hectare} queimados (acréscimo de \(\approx 400\%\) em relação aos dois anos anteriores), com emissão estimada de \SI{791}{\mega\tonne} de \ce{CO2}. Esses números reforçam a urgência operacional do desenvolvimento de ferramentas modernas de previsão e vigilância, capazes de produzir mapas de risco em tempo quase-real e auditáveis por gestores ambientais.

% ============================================================================
% (B) NOVA REDACAO DOS OBJETIVOS ESPECIFICOS -- substituir o conteudo do
%     ambiente que lista os 6 objetivos especificos (Secao 1.2.2).
% ============================================================================

\begin{enumerate}
  \item Realizar uma revisão bibliográfica sobre as técnicas de aprendizado de máquina e classificação em evidência na literatura, incluindo trabalhos publicados entre 2024 e 2025 que utilizam \emph{ensembles} e métodos de interpretabilidade aplicados a risco de incêndio florestal.

  \item Mapear as regiões da Amazônia Legal que sofreram com queimadas durante a última década (2014--2023) e definir suas principais causas, com base em registros do Programa Queimadas/INPE.

  \item Obter e enriquecer os dados climáticos, topológicos e de incêndios históricos, incluindo enriquecimento com NASA POWER (umidade relativa RH2M) e NASA FIRMS (FRP em tempo quase-real), além da construção de \textbf{24 \emph{features} físico-climáticas avançadas} (Tier~1: SPI, KBDI \emph{proxy}, VPD \emph{proxy}, anomalias de precipitação, histórico estendido de fogo) embasadas na literatura recente~\cite{seager2015climatology,forests2024cerradoamazon,npjhazards2025fire}.

  \item Desenvolver um algoritmo que utilize diferentes modelos de aprendizado de máquina para a previsão de regiões mais suscetíveis a queimadas, comparando oito abordagens (Regressão Logística, SGD, Random Forest com SMOTE, RF balanceado otimizado por \textbf{Optuna}~\cite{akiba2019optuna}, XGBoost~\cite{chen2016xgboost}, LightGBM~\cite{ke2017lightgbm}, CatBoost~\cite{prokhorenkova2018catboost}, e \emph{ensembles} \emph{Voting Soft} e \emph{Stacking}~\cite{wolpert1992stacked}), com calibração isotônica de probabilidades~\cite{zadrozny2002transforming} e ajuste multi-classe de limiares de decisão.

  \item Avaliar a eficácia dos modelos propostos por meio de métricas de desempenho por classe (precisão, revocação, \(\text{F}_1\), \(\text{F}_1\)-macro), validação cruzada estratificada \(k\)-\emph{fold} e \textbf{validação temporal estrita} (\emph{rolling-origin} por ano)~\cite{bergmeir2012cv} para quantificar o viés do \emph{split} aleatório, complementadas por análise de interpretabilidade via \textbf{SHAP}~\cite{lundberg2017unified,lundberg2020local} e \textbf{\emph{Permutation Importance}}~\cite{breiman2001random,fisher2019all}.

  \item Desenvolver uma \textbf{aplicação web interativa explicável} (\emph{Flask} + \emph{Folium}~\cite{folium,leafletjs}) que, ao receber um clique do usuário sobre qualquer ponto da Amazônia Legal, classifique o risco do ponto consultado e exiba a explicação local das \emph{features} responsáveis pela decisão, viabilizando uso operacional por gestores ambientais.
\end{enumerate}

% ============================================================================
% (C) NOVA REDACAO DA ESTRUTURA DA MONOGRAFIA -- substituir o conteudo da
%     Secao 1.3 (paginas 16--17 do original) por este bloco.
% ============================================================================

Este trabalho está organizado em seis capítulos que abordam diferentes etapas da pesquisa, proporcionando uma compreensão abrangente do desenvolvimento do sistema de previsão de queimadas na Amazônia.

\paragraph{Capítulo 2 -- Estado da Arte.} São explorados trabalhos correlatos e avanços significativos na área de previsão de queimadas, com foco em estudos que propõem o mapeamento de regiões suscetíveis a fogo. A revisão aborda a extensão do problema considerando seus impactos ambientais, sociais e econômicos. Adicionalmente, são examinadas abordagens tradicionais de aprendizado de máquina e o início dos sistemas de prevenção, complementadas por uma seção de \textbf{literatura recente 2024--2025} (\emph{ensembles} explicáveis, índice KBDI no Cerrado-Amazônia, validação temporal estrita).

\paragraph{Capítulo 3 -- Fundamentação Teórica.} Estabelece as bases conceituais necessárias para a compreensão do problema das queimadas na Amazônia. Inicia-se com a contextualização do bioma amazônico e a problemática dos incêndios florestais, seguida pela discussão do sensoriamento remoto e das tecnologias espaciais. O capítulo prossegue com os fundamentos da inteligência artificial e do aprendizado de máquina, contemplando \textbf{modelos clássicos} (Regressão Logística, Random Forest), \textbf{modelos modernos de \emph{boosting}} (XGBoost, LightGBM, CatBoost), \textbf{\emph{ensembles}} (Voting e Stacking), \textbf{otimização de hiperparâmetros} (Optuna), \textbf{interpretabilidade} (SHAP e Permutation Importance), \textbf{tratamento de desbalanceamento}, \textbf{calibração de probabilidades}, \textbf{índices físico-climáticos de fogo} (KBDI, SPI, VPD, De Martonne) e \textbf{validação temporal} de modelos.

\paragraph{Capítulo 4 -- Metodologia.} Descreve detalhadamente os passos e procedimentos adotados para desenvolver o sistema de previsão. São delineados os métodos de coleta e enriquecimento dos dados (Programa Queimadas + NASA POWER + NASA FIRMS), o pré-processamento, a engenharia das 24 \emph{features} Tier~1, o protocolo de treinamento dos modelos, a avaliação (métricas, interpretabilidade, \emph{threshold tuning}, validação temporal), e o desenvolvimento da aplicação web interativa explicável.

\paragraph{Capítulo 5 -- Resultados.} Apresenta os resultados experimentais: comparativo entre as oito abordagens treinadas, desempenho detalhado do modelo final, análise de interpretabilidade por SHAP e \emph{Permutation Importance}, resultados da validação temporal estrita, efeito do \emph{threshold tuning}, evolução incremental do desempenho ao longo das decisões metodológicas e validação operacional via consultas comparativas no aplicativo web.

\paragraph{Capítulo 6 -- Conclusão.} Sintetiza os principais achados, registra as limitações reconhecidas e enumera os trabalhos futuros prioritários para evoluir o sistema rumo à integração de novas fontes (NDVI/EVI MODIS, MapBiomas Fogo) e de causalidade estrita nas \emph{features} de janela móvel.
